\documentclass{article}
\usepackage{amsmath}
\usepackage{amssymb}
\usepackage{fullpage}
\usepackage{enumerate}
\usepackage{xcolor}
\pagestyle{empty}
\usepackage{graphicx}
\usepackage{bm}
\graphicspath{{C:\Users\steve\Pictures}}
\begin{document}

\noindent{{\Large \bf Fall 25, 4100 Problem Set 6, Hayden Fuller}

\large

\bigskip
\noindent{Please complete all problems. For any proof questions, please structure your proof similar to those from lecture.

\begin{enumerate}

\item Given $A=\left[
\begin{array}{cccc}
2 & -4 & 1  & 3\\
3 &  -6  & 0 & -3\\
-1 &  2 & 1 &  0\\
4 &  -8& -1  & 3
\end{array}\right]$
compute adj$(A)$ and use it to find bases for ${\rm null}(A)$ and ${\rm null}(A^T)$
\\
\\$\textrm{adj}(A)=\left[
\begin{array}{cccc}
-36 & 0 & 72  & 36\\
-18 &  0  & 36 & 18\\
0 &  0 & 0 &  0\\
0 &  0 & 0  & 0
\end{array}\right]$
\\$\textrm{det}(A)=0$, $\textrm{rank}(A)=3$, $\textrm{dim ker}(A)=1$, $\textrm{dim ker}(A^T)=1$
\\basis vectors:
\\$\textrm{ker}(A)=\textrm{span}\{(2,1,0,0)^T\}$, $\textrm{ker}(A^T)=\textrm{span}\{(-1,0,2,1)^T\}$
\\$A(2,1,0,0)^T=0$, $A^T(-1,0,2,1)^T=0$ 
\\$\textrm{adj}(A)=18\left[\begin{array}{c}2\\1\\0\\0\end{array}\right]\left[\begin{array}{cccc}-1&0&2&1\end{array}\right]$


\newpage\item Find the eigenvalues and eigenvectors given the following matrices:
\begin{enumerate}[a.]
\item $A= \left[
\begin{array}{cc}
4 & 3 \\
5 & 6 \\
\end{array} \right] $
\\
\\$\textrm{det}(A-\lambda I)=(4-\lambda)(6-\lambda)-15=\lambda^2-10\lambda+9$
\\$\lambda=9$, $\lambda=1$
\\$(A-9I)v=0$, $v=\left[\begin{array}{c}3\\5\end{array}\right]$
\\$(A-1I)v=0$, $v=\left[\begin{array}{c}1\\-1\end{array}\right]$
\\
\item $B= \left[
\begin{array}{ccc}
1 & 2 & -1 \\
0 & 0 & 4 \\
0 & 2 & -2
\end{array} \right] $
\\
\\$\textrm{det}(B-\lambda I)=(1-\lambda)(\lambda^2+2\lambda-8)=(1-\lambda)(\lambda-2)(\lambda+4)$
\\$\lambda=1$, $\lambda=2$, $\lambda=-4$
\\$(B-1I)v=0$, $v=\left[\begin{array}{c}1\\0\\0\end{array}\right]$
\\$(B-2I)v=0$, $v=\left[\begin{array}{c}3\\2\\1\end{array}\right]$
\\$(B+4I)v=0$, $v=\left[\begin{array}{c}3\\-5\\5\end{array}\right]$
\end{enumerate}


\newpage\item Let $T\in\mathcal{L}(\mathbb{C}^n,\mathbb{C}^n)$ be such that for linear independent $\bm{v}_1,\bm{v}_2\in\mathbb{C}^n$ we have that $T\bm{v}_1=\bm{v}_1+6\bm{v}_2$ and $T\bm{v}_2=3\bm{v}_1+4\bm{v}_2$. Find two eigenvalues and associated eigenvectors of $T$.
\\
\\$M=\left[\begin{array}{cc}1&3\\6&4\end{array}\right]$
\\$\textrm{det}(M-\lambda I)=(1-\lambda)(4-\lambda)-18=\lambda^2-5\lambda-14=0$
\\$\lambda=7$, $\lambda=-2$
\\$(M-7I)v=0$, $v=\left[\begin{array}{c}1\\2\end{array}\right]$
\\$1v_1+2v_2$, $7$
\\$(M+2I)v=0$, $v=\left[\begin{array}{c}1\\-1\end{array}\right]$
\\$1v_1-1v_2$, $-2$
\\$T(v_1+2v_2)=7(v_1+2v_2)$
\\$T(v_1-v_2)=-2(v_1-v_2)$


\newpage\item Let $A$ and $B$ be similar such that $A=M^{-1}BM$. Prove that $A$ and $B$ have the same characteristic and minimal polynomial. 
\\
\\let $A=M^{-1}BM$
\\then $\chi_A(\lambda)=\textrm{det}(\lambda I-A)=\textrm{det}(\lambda I-M^{-1}BM)=\textrm{det}(M^{-1}(\lambda I-B)M)$
\\$\textrm{det}(M^{-1}(\lambda I-B)M)=\textrm{det}(M^{-1})\textrm{det}(\lambda I-B)\textrm{det}(M)=\textrm{det}(\lambda I-B)=\chi_B(\lambda)$
\\thus $\chi_A(\lambda)=\chi_B(\lambda)$
\\
\\let $m_B(\lambda)$ be the minimal polynomial of $B$
\\consider $m_B(A)$
\\$m_B(A)=m_B(M^{-1}BM=M^{-1}m_B(B)M$
\\but $m_B(B)=0$, so $m_B(A)=M^{-1}0M=0$
\\hence $m_B$ is also an annihilating polynomial for $A$, so the minimal polynomial $m_A$ of $A$ divides $m_B$
\\by symmetry, $m_B$ of $B$ divides $m_A$
\\two monic polynomails that divide each other must be equal, so $m_A=m_B$


\newpage\item Find the set of all $A\in\mathbb{R}^{2\times 2}$ such that $A^6=I$ up to similarity classes by finding all possible minimal polynomials. 
\\
\\$m_A(x)$ divides $x^6-1$
\\$x^6-1=(x^3-1)(x^3+1)=(x-1)(x^2+x+1)(x+1)(x^2-x+1$
\\possible minimal polynomials $m(x)=x-1,x+1,x^2+x+1,x^2-x+1,x^2-1$
\\$m(x)=x-1$, then $A=I$
\\$m(x)=x+1$, then $A=-I$
\\$m(x)=x^2-1=(x-1)(x+1)$, then $A$ is diagonalizable with eigenvalues 1 and -1
\\$A\sim\left[\begin{array}{cc}1&0\\0&-1\end{array}\right]$
\\$m(x)=x^2+x+1$, then 
\\$A\sim\left[\begin{array}{cc}-\frac{1}{2}&-\frac{\sqrt{3}}{2}\\\frac{\sqrt{3}}{2}&-\frac{1}{2}\end{array}\right]$, satisfying $A^3=I$ and $A\neq I$, so $A^6=I$
\\$m(x)=x^2-x+1$, then $A\sim\left[\begin{array}{cc}\frac{1}{2}&-\frac{\sqrt{3}}{2}\\\frac{\sqrt{3}}{2}&\frac{1}{2}\end{array}\right]$, satisfying $A^6=I$ and $A^3\neq I$
\\$A\in\{I,-I,\left[\begin{array}{cc}1&0\\0&-1\end{array}\right],\left[\begin{array}{cc}-\frac{1}{2}&-\frac{\sqrt{3}}{2}\\\frac{\sqrt{3}}{2}&-\frac{1}{2}\end{array}\right],\left[\begin{array}{cc}\frac{1}{2}&-\frac{\sqrt{3}}{2}\\\frac{\sqrt{3}}{2}&\frac{1}{2}\end{array}\right]\}$

\newpage\item Let $A\in\mathbb{R}^{n\times n}$ be invertible and have characteristic polynomial of $p(\lambda)$. Show that $A^{-1}$ will have characteristic polynomial of 
$$q(\lambda)=\displaystyle{\frac{(-\lambda)^n}{|A|}~p\left(\frac{1}{\lambda}\right)}$$
\\
\\$p(\frac{1}{\lambda})=\textrm{det}(\frac{1}{\lambda}I-A)=\textrm{det}(\frac{1}{\lambda}(I-\lambda A))=(\frac{1}{\lambda})^n\textrm{det}(I-\lambda A)$
\\since $\textrm{det}(I-\lambda A)=(-1)^n\textrm{det}(\lambda A-I)$,
\\this gives $p(\frac{1}{\lambda})=\frac{(-1)^n}{\lambda^n}\textrm{det}(\lambda A-I)$
\\hence $\textrm{det}(\lambda A-I)=(-\lambda)^np(\frac{1}{\lambda})$
\\$\textrm{det}(\lambda I-A^{-1})=\textrm{det}(A^{-1})\textrm{det}(\lambda A-I)=\frac{1}{\textrm{det}(A)}(-\lambda)^np(\frac{1}{\lambda})$
\\therefore $\textrm{det}(\lambda I-A^{-1})=\frac{(-\lambda)^n}{\textrm{det}(A)}p(\frac{1}{\lambda})$
\\QED

\end{enumerate}

\end{document}
